\documentclass{umk-matematika}

\begin{document}

\maketitlepage
    {Практическая работа №4}
    {Решение тригонометрических уравнений}
    {}

\section{Фундамент (1--4 балла)}

\textbf{Задача 1 (1 балл).} Решите уравнение: sin x = 0

\textbf{Задача 2 (2 балла).} Решите уравнение: cos x = 1

\textbf{Задача 3 (3 балла).} Решите уравнение: sin x = 12

\textbf{Задача 4 (4 балла).} Решите уравнение: cos x = √22

\section{Стены (5--7 баллов)}

\textbf{Задача 5 (5 баллов).} Решите уравнение: 2 sin x – 1 = 0

\textbf{Задача 6 (6 баллов).} Решите уравнение: cos x + 1 = 0

\textbf{Задача 7 (7 баллов).} Решите уравнение: tg x – 1 = 0

\section{Кровля (8--10 баллов)}

\textbf{Задача 8 (8 баллов).} Решите уравнение: 2 cos² x – 1 = 0

\textbf{Задача 9 (9 баллов).} Решите уравнение на [0; 2π]: sin x = –12

\textbf{Задача 10 (10 баллов).} Угол наклона кровли α удовлетворяет уравнению tg α = 0,5. Найдите α в градусах (округлите до целых).

\newpage
\section*{Ответы}

\begin{enumerate}
  \item x = πk, k ∈ Z.
  \item x = 2πk, k ∈ Z.
  \item x = (–1)k · π6 + πk, k ∈ Z.
  \item x = ± π4 + 2πk, k ∈ Z.
  \item x = (–1)k · π6 + πk, k ∈ Z.
  \item x = π + 2πk, k ∈ Z.
  \item x = π4 + πk, k ∈ Z.
  \item x = π4 + πk2, k ∈ Z.
  \item 7π6; 11π6.
  \item α ≈ 27$^\circ$.
\end{enumerate}

\end{document}